% ═══════════════════════════════════════════════════════════════════════════
% LÖSUNGEN DS-04: Evaluación 1 (Selbstkontrolle Gruppe A)
% Basiert auf A_tope.com Evaluación Unidad 1
% ═══════════════════════════════════════════════════════════════════════════

\documentclass[11pt, a4paper]{article}

\usepackage[utf8]{inputenc}
\usepackage[T1]{fontenc}
\usepackage[ngerman]{babel}
\usepackage{helvet}
\renewcommand{\familydefault}{\sfdefault}

\usepackage[margin=1.5cm, top=1.2cm, bottom=1.2cm]{geometry}
\usepackage[table]{xcolor}
\usepackage{array}
\usepackage{tabularx}
\usepackage{booktabs}
\usepackage{tikz}
\usepackage{tcolorbox}
\usepackage{enumitem}
\usepackage{amssymb}
\tcbuselibrary{skins}

% Farben
\definecolor{spanisch}{HTML}{c41e3a}
\definecolor{gruppeA}{HTML}{2a7a4b}
\definecolor{hellgrau}{HTML}{f5f5f5}
\definecolor{loesung}{HTML}{c41e3a}

\pagestyle{empty}

% Lösungs-Formatierung
\newcommand{\lsg}[1]{\textcolor{loesung}{\textbf{#1}}}

\begin{document}

% ═══════════════════════════════════════════════════════════════════════════
% TITEL
% ═══════════════════════════════════════════════════════════════════════════

\begin{center}
\begin{tikzpicture}
\fill[gruppeA] (0,0) rectangle (18,0.9);
\node[white, font=\Large\bfseries] at (9,0.45) {Evaluación 1 -- Lösungen zur Selbstkontrolle};
\end{tikzpicture}
\end{center}

\vspace{0.2cm}

\begin{center}
\textit{Überprüfe deine Antworten und markiere, was du noch üben solltest!}
\end{center}

\vspace{0.3cm}

% ═══════════════════════════════════════════════════════════════════════════
% AUFGABE 1: SER / ESTAR / HAY
% ═══════════════════════════════════════════════════════════════════════════

\begin{tcolorbox}[
    colback=hellgrau,
    colframe=spanisch,
    title={\color{white}\bfseries Aufgabe 1: Completa con ser, estar o hay},
    fonttitle=\bfseries
]

\begin{enumerate}[label=\alph*)]
    \item En mi barrio \lsg{hay} un parque grande.
    \item Mi hermana \lsg{es} profesora.
    \item ¿Dónde \lsg{está} el supermercado?
    \item Los libros \lsg{están} en la mesa.
    \item Madrid \lsg{es} la capital de España.
    \item \lsg{Hay} muchas tiendas en el centro.
    \item ¿Cómo \lsg{estás}? -- \lsg{Estoy} bien, gracias.
    \item La biblioteca \lsg{está} al lado del banco.
    \item Mis padres \lsg{son} de Alemania.
    \item Hoy no \lsg{hay} clase.
\end{enumerate}

\vspace{0.2cm}
\footnotesize\textbf{Tipp:} hay = es gibt (Existenz) | estar = wo/wie (Ort/Zustand) | ser = was/wer (Identität)
\end{tcolorbox}

\vspace{0.3cm}

% ═══════════════════════════════════════════════════════════════════════════
% AUFGABE 2: KONJUGATION IR
% ═══════════════════════════════════════════════════════════════════════════

\begin{tcolorbox}[
    colback=hellgrau,
    colframe=spanisch,
    title={\color{white}\bfseries Aufgabe 2: Conjuga el verbo IR},
    fonttitle=\bfseries
]

\begin{tabular}{@{}p{6cm}p{10cm}@{}}
\begin{enumerate}[label=\alph*)]
    \item yo \lsg{voy}
    \item tú \lsg{vas}
    \item él/ella \lsg{va}
    \item nosotros \lsg{vamos}
    \item vosotros \lsg{vais}
    \item ellos \lsg{van}
\end{enumerate}
&
\textbf{Ergänze die Sätze:}
\begin{enumerate}[label=\alph*), start=7]
    \item Yo \lsg{voy} al colegio en autobús.
    \item ¿Adónde \lsg{vas} tú?
    \item Nosotros \lsg{vamos} a la playa mañana.
    \item María \lsg{va} a estudiar medicina.
\end{enumerate}
\end{tabular}

\vspace{0.2cm}
\footnotesize\textbf{Tipp:} ir + a + Ort | ir + a + Infinitiv (Zukunft) | ir + en + Transportmittel
\end{tcolorbox}

\vspace{0.3cm}

% ═══════════════════════════════════════════════════════════════════════════
% AUFGABE 3: MODALVERBEN
% ═══════════════════════════════════════════════════════════════════════════

\begin{tcolorbox}[
    colback=hellgrau,
    colframe=spanisch,
    title={\color{white}\bfseries Aufgabe 3: Completa con poder, querer o tener que},
    fonttitle=\bfseries
]

\begin{enumerate}[label=\alph*)]
    \item Yo \lsg{quiero} ser médico. (wollen)
    \item ¿\lsg{Puedes} ayudarme? (können)
    \item Mañana \lsg{tengo que} trabajar. (müssen)
    \item Nosotros \lsg{queremos} viajar a España. (wollen)
    \item Ella no \lsg{puede} venir hoy. (können)
    \item Los alumnos \lsg{tienen que} hacer los deberes. (müssen)
\end{enumerate}

\vspace{0.2cm}
\footnotesize\textbf{Achtung:} tener que = müssen (mit \textbf{que}!) | poder = können (o→ue) | querer = wollen (e→ie)
\end{tcolorbox}

\vspace{0.3cm}

% ═══════════════════════════════════════════════════════════════════════════
% AUFGABE 4: WORTSCHATZ
% ═══════════════════════════════════════════════════════════════════════════

\begin{tcolorbox}[
    colback=hellgrau,
    colframe=spanisch,
    title={\color{white}\bfseries Aufgabe 4: Vocabulario -- Traduce al español},
    fonttitle=\bfseries
]

\begin{tabular}{@{}p{8cm}p{8cm}@{}}
\begin{enumerate}[label=\alph*)]
    \item das Stadtviertel: \lsg{el barrio}
    \item die Apotheke: \lsg{la farmacia}
    \item die Bushaltestelle: \lsg{la parada de autobús}
    \item neben: \lsg{al lado de}
    \item gegenüber: \lsg{enfrente de}
\end{enumerate}
&
\begin{enumerate}[label=\alph*), start=6]
    \item ruhig: \lsg{tranquilo/a}
    \item modern: \lsg{moderno/a}
    \item nah: \lsg{cerca (de)}
    \item die Bäckerei: \lsg{la panadería}
    \item das Einkaufszentrum: \lsg{el centro comercial}
\end{enumerate}
\end{tabular}
\end{tcolorbox}

\vspace{0.4cm}

% ═══════════════════════════════════════════════════════════════════════════
% SELBSTEINSCHÄTZUNG
% ═══════════════════════════════════════════════════════════════════════════

\begin{center}
\fcolorbox{spanisch}{hellgrau}{%
\begin{minipage}{0.95\textwidth}
\textbf{Selbsteinschätzung:} Zähle deine korrekten Antworten!

\vspace{0.2cm}

\begin{tabular}{|l|c|c|c|c|}
\hline
\rowcolor{spanisch!20}
\textbf{Aufgabe} & \textbf{Maximum} & \textbf{Meine Punkte} & \textbf{$\geq$ 70\%} & \textbf{Üben?} \\
\hline
1. ser/estar/hay & 10 & \_\_\_ / 10 & $\geq$ 7 & $\square$ \\
\hline
2. ir (Konjugation) & 10 & \_\_\_ / 10 & $\geq$ 7 & $\square$ \\
\hline
3. Modalverben & 6 & \_\_\_ / 6 & $\geq$ 4 & $\square$ \\
\hline
4. Wortschatz & 10 & \_\_\_ / 10 & $\geq$ 7 & $\square$ \\
\hline
\rowcolor{gruppeA!20}
\textbf{Gesamt} & \textbf{36} & \_\_\_ / 36 & $\geq$ 25 & \\
\hline
\end{tabular}

\vspace{0.2cm}

\footnotesize Kreuze an, was du noch üben solltest. Frag E.H. oder die Lehrkraft bei Unklarheiten!
\end{minipage}%
}
\end{center}

\end{document}
